%!TEX program = lualatex
\documentclass[ngerman,headheight=70pt]{scrartcl}
\usepackage[ngerman]{babel}
\usepackage{blindtext}
\usepackage[utf8]{luainputenc}
\usepackage{scrlayer-scrpage}
\usepackage[T1]{fontenc}
\usepackage{fontspec}
\usepackage{multicol}
\usepackage{enumitem}
\usepackage[german=quotes]{csquotes}
\usepackage[bottom=2cm,footskip=8mm]{geometry}
\setmainfont{Source Sans Pro}
\MakeOuterQuote{"}

\setlength{\parindent}{0pt}
\setlength{\parskip}{6pt}

\pagestyle{plain.scrheadings}
\thispagestyle{scrheadings}
\lohead{\vspace*{-1.5cm}\includegraphics[scale=0.4]{../images/stupa_logo_rot.pdf}}
\rohead{\\Präsidium des\\Studierendenparlaments}

\lofoot[Universität Hamburg · Präsidium des Studierendenparlaments\\\hrule Von-Melle-Park 5 · D-20146 Hamburg]{Universität Hamburg · Präsidium des Studierendenparlaments\\\hrule Von-Melle-Park 5 · D-20146 Hamburg}
\cofoot[]{}
\cohead[\thepage]{}

\renewcommand{\thesubsection}{TOP \arabic{subsection}}
\renewcommand{\thesubsubsection}{\arabic{subsubsection}.}
\setcounter{subsection}{-1}

\begin{document}
    UHH · StuPa-Präsidium · Von-Melle-Park 5 · D-20146 Hamburg

    %%% PAGE ONE (EDIT LATER) %%%%
    % Hier den Titel des Protokolls eintragen
    \section*{Protokoll der 13. Sitzung des Studierendenparlaments vom 12. Januar 2017}

    % Protokollanten eintragen
    \textbf{Protokoll: Jim Martens}\\
    \textbf{Ort: VMP5 0077}\\
    \textbf{Beginn: 20.54 Uhr}\\
    \textbf{Ende: 23.55 Uhr}

    % Wohl die komplizierteste Stelle im Protokoll
    % Als Faustregel sollten nicht mehr als 3 Namen auf einer Zeile stehen
    % Einfach an dieser Vorlage orientieren
    % Im Idealfall ist hier nicht viel zu ändern, da kaum jemand fehlt.
    % Rücktritte müssen natürlich pro Protokoll angepasst werden
    \vspace{0.5cm}
    \begin{tabular}{ll}
        Anwesend: & \\
            RCDS (5 Sitze): & Ramon Weilinger \\
             CampusGrün (14 Sitze): & Laura Franzen, Geoffrey Youett, Philipp Droll, \\
                                   & Yasemin Günther, Christian Lagod, Jim Martens, \\
                                   & Svenja Horn, Martin Sievert, Armin Günther \\
             BGZM-Liste (2 Sitze) : & \\
             WiWi (2 Sitze): & \\
             Unicorns (5 Sitze): &    \\
             Liste LINKS (3 Sitze): & Gunhild Berdal, Till Petersen, Sinah Mielich \\
             HWP (2 Sitze): &  \\
             MIN (4 Sitze): &  \\
             SDS* (3 Sitze): & Mena Winkler, Jacob Petersein, Artur Brückmann \\
             Bart-LISTE (2 Sitze): & Timo Zeimet, Dominic Laumer \\
             LHG (1 Sitz): & Tobias Heisig \\
             harte zeiten (1 Sitz): & Tobias Berking \\
             Jura (1 Sitz): & \\
             AL (2 Sitze): &  \\
            & \\
        Entschuldigt: & Antonia Niecke, Ramin Shakiba, Jennifer Maack, \\
                      & Benjamin Welling (RCDS), Mario Moldenhauer, Martin Sievert (CG) \\
                      & Jakob Pape, Karen Martirosian (BGZM), Claas-Friso Hente, \\
                      & Charleen Lorenz (WiWi), Katharina Kucza, Johannes Peplow, \\
                      & Annkathrin Löffler, Andreas Hartkamp (Unicorns), Hakan Tasdemir, \\
                      & Gökce Celik (HWP), Ailina Salten, Lotte Rullkötter, \\
                      & Nabil Jabareen, Jan Detampel (MIN), Karima Schulze\\
                      & Henri Weber (AL) \\
                      &\\
        Unentschuldigt abwesend: & Elena Rysikova, Melf Johannsen, Ines Friedrich \\
                                & Dean Haltmeyer (CG), Marielle Hermstrüwer (Unicorns), \\
                                & Maurice Martens (Jura)\\
        Rücktritte: & \\
        Mandatsverluste: & \\
    \end{tabular}
    \newpage
    %%% Vorgeschlagene Tagesordnung %%%
    \underline{Vorgeschlagene Tagesordnung}
    \begin{enumerate}[label={\textbf{Top \theenumi}},leftmargin=*]
        \item Beitragsordnung (V1617-046) (30 Min.)
        \item VG-Wort (V1617-047) (30 Min.)
        \item Verschiedenes (5 Min.)
    \end{enumerate}

    \newpage

    % Hier beginnt Teil A. Bei Fortsetzungen kommt hier nur der Teil über
    % die Feststellung der endgültigen Fassung der Tagesordnung.
    % Für sonstige Sitzungen ist dies bereits das Gerüst für die Formalia.
    % Man muss nur noch die protokollierten Inhalte eintragen.

    {\Large \underline{Teil A}}

    \subsection{Formalia}

    \textit{
    Gunhild eröffnet die Sitzung um 20.54 Uhr. Ramin vertritt Ramon im Präsidium.
    }

    \subsubsection{Geschäftsbericht Präsidium}

    Gunhild berichtet aus dem Präsidium. Beschäftigung mit der Wahlvorbereitung.
    Wahlhelfer*innen haben viele Dinge verteilt - auch an Außenstandorte und
    Wohnheime. Letzte Woche wurden die fristgerechten Briefe von der Briefwahl
    gezählt. <Zahlen einfügen> Starke Beteiligung bei der Urnenwahl wichtig.
    Es ist wichtig Leute zu informieren, dass sie an der Urne überprüfen, ob
    ihre Briefwahlstimme angekommen ist.

    Urnen werden morgen zusammen mit Drexler versiegelt. Montag geht es mit der
    Urnenwahl los.

    Sonderproblem mit dem Studienkolleg. Infomaterial wurde dort hingebracht
    und am Freitag, 14.1., findet ein Infotreff dort statt.

    Gunhild ruft zu Besetzung von Urnen auf. Wenn Leute in FSRen aktiv sind, sollte
    dort für die Wahl mobilisiert werden. Nehmt auch an der Wahl zum AS statt.

    Ungefähr 10 Briefumschläge wurden in Briefen gefunden. AS-Briefwahlumschläge
    können auch an den Urnen abgegeben werden. Am Dienstag, 17.1., werden die
    Briefe um 13 Uhr zum Wahlamt gebracht.

    Jim verliest die Rücktritte. Vorschlag Rücktrittsanteil in Wahlzeitung
    aufzunehmen.

    \subsubsection{Anfragen an das Präsidium}

    Philipp merkt an, dass viele Plakate zerstört wurden und fragt, ob auch andere
    Gruppen betroffen sind.

    Karima merkt an, dass auch bei ihnen Plakate kaputt gegangen sind. Sie vermutet
    jedoch eher eine betrunkene Bande.

    Tobias hat auch einige Verluste zu beklagen, aber nichts systematisches.

    Laura merkt an, dass z.B. um das Geomatikum alle Plakate restlos verschwunden
    sind.

    Melf weist auf Orte hin, wo Hausmeister Plakate wegräumen. Um das Geomatikum
    ist das jedoch eher nicht der Fall.

    Paula merkt an, dass Richtung Dammtor und Hauptcampus alles um Silvester
    weg war.

    Karima merkt an, dass man aufpassen sollte, ob Plakate von gewissen Listen
    gezielt entfernt werden.

    Gunhild bietet an, dass der Beschluss "Fairer Wahlkampf" stärker auf der
    Website platziert werden kann. Die Hausmeister könnten wir informieren.

    Franzi weist darauf hin, dass die Hausmeister der HWP über die bestellten
    Tische informiert haben. Gunhild klärt auf, dass die Tische nicht vom
    Präsidium zu den Gebäuden geschleppt werden müssen, sondern direkt zu den
    Gebäuden geliefert werden.

    \subsubsection{Geschäftsbericht AStA}

    21:15

    Franzi:
    Dies im WS 17/18 wird über Finanzen stattfinden. Der AStA soll der Uni sagen,
    mit welchen Leuten diese Vorbereitung zum Dies stattfinden soll. Die Gruppe
    wollte sich bereits im Januar konstituieren.

    Vertrag mit Sprickenhof AG und Behörde für Wissenschaft, Forschung und
    Gleichstellung unterzeichnet:
    Überseering 35 wird die Interimsunterbringung für den Philturm sein. Bautermin
    bzgl. Philturm Sommer 2017 steht. Bibliotheken und gesamter Philturm kann
    in das Interimsgebäude ziehen. Herausforderung für studentische und politische
    Kultur mit großer Entfernung zum Hauptcampus.

    Kantine im Gebäude wird privatwirtschaftlich betrieben. Dem Studierendenwerk
    entfällt die Philmensa. Im Januar wird es Vollversammlung des Philturms
    unter Leitung der Uni geben.

    RPO kontrovers im AS diskutiert. "Rechter Pol" musste beteuern, dass sie
    nichts gegen Studienreform haben. Ritter und Rupp haben versprochen einen Entwurf
    für eine General Studies Umsetzung bis zur RPO-Tagung zur Verfügung zu stellen.

    AS hat beschlossen, dass in 2. Sitzung der neuen Legislatur erneut eine 1.
    Lesung eines RPO-Entwurfs stattfindet. Dies zur Studienreform im Mai gibt
    es weiterhin. Lenzen hat betont, dass eine RPO notwendig ist und dass es sie
    geben wird.

    Armin:
    Am Dienstag hat die 3. Sitzung der G20-Plattform stattgefunden. Es wurde
    in G20-Plattform konkrete Arbeitsweise entwickelt.

    Demo in April, bis dahin Mobilisierung.

    Information über Chancen eines Volksentscheids. Am 24.1. um 19 Uhr nächstes
    Treffen.

    Artur:
    Hinweis auf Veranstaltung namens "Putsch im Putsch in der Türkei und
    Berufsverbote in Deutschland" (Name des Gastes ergänzen)
    Datum: 21. Januar 19 Uhr

    Geoffrey:
    Das neue Stille Wörtchen ist fertig mit Schwerpunkt Wahlen, Rückblick, AS
    und was so gewählt wird. Anwesende von Außenstandorten sind eingeladen das
    StiWö mitzunehmen und aufzuhängen.

    \subsubsection{Anfragen an den AStA}

    @Armin
    Tobias ist nicht klar geworden, wie die Gesellschaft überzeugt werden soll.
    Welche Bündnispartner sind angedacht? Wie wird mit DGB-Entscheidung umgegangen?

    Armin:
    Manifest ist guter Start, um es in die Gesellschaft zu tragen.
    AStA-Infobroschüre zum G20 kann auch hilfreich dahingehend sein. Langer Weg
    bis gesamtgesellschaftlichem Widerstand. Mitglieder in DGB und Parteien
    motivieren gegen Vorstandsentscheidung anzugehen.

    Franzi:
    Es wird in irgendeiner Form eine Unterschriftenkampagne geben.

    \subsubsection{Dringlichkeitsanträge des AStA}

    Keine.

    \subsubsection{Feststellung der endgültigen Fassung des Teils B der Tagesordnung}

    \newpage
    %%% BESCHLOSSENE TAGESORDNUNG %%%
    \underline{Endgültige Fassung der Tagesordnung}
    \begin{enumerate}[label={\textbf{Top \theenumi}},leftmargin=*]
        \item Beitragsordnung (V1617-046) (30 Min.)
        \item Manifest zu G20 (V1617-049) (45 Min.)
        \item VG-Wort (V1617-047) (30 Min.)
        \item Englischsprachige Wahlwerbung (V1617-048) (30 Min.)
        \item Verschiedenes (5 Min.)
    \end{enumerate}

    Mehrheit zu 0 zu 2 angenommen.

    \subsubsection{Feststellung der Beschlussfähigkeit}

    \textit{Mit 28 Parlamentarier*innen beschlussfähig.}

    \subsubsection{Genehmigung der Protokolle der vorangegangenen Sitzungen}

    Keine Protolle zu genehmigen.

    \vspace{0.5cm}
    {\Large \underline{Teil B}}

    21.43

    \subsection{Beitragsordnung}

    Philipp stellt den Antrag vor.

    Franzi stellt vor, dass 0,10€ Verringerung des AStA-Budgets eine Veränderung
    von ca. 8400€ bedeutet.

    Tobias fragt, ob der AStA sich überlegt hat, die Forderungen des HVV nicht
    zu bedienen. Philipp sagt, dass der AStA frühzeitiger sich engagiert hat,
    dies aber nicht gefruchtet hat.

    Gunhild - frühere Asten, offensive in der Frage

    Artur merkt an, dass nach längerer Diskussion die Vertreter des HVV die
    Punkte des AStA im HVV intern diskutieren wollten. Bei einer Busbeschleunigungsfahrt
    wurde mit dem zuständigen Staatsrat gesprochen und der hat sich gegen die Verteilung
    des Senatsprozents an die Studierenden ausgesprochen, weil dann ja alle kommen
    könnten

    Jennifer bittet, dass in der Debatte um Beitragsordnung sich nicht über
    kostenlosen ÖPNV unterhalten wird.

    Tobias findet, dass man sich zusammen mit den anderen ASten ab sofort
    zusammenschließen sollte, um eine Mobilisierung gegen die Erhöhungen zu
    erreichen. Tobias fragt, warum in der Vergangenheit nichts dahingegend passiert
    ist.

    Jennifer fragt, warum denn kein entsprechender Antrag gestellt wurde.

    Till erinnert an aktivere Zeiten. Es muss durch Öffentlichkeitsarbeit und
    Skandalisierung kontinuierlich auf das Problem hingewiesen wird. Die Botschaft
    sollte lauten: Semesterticket ja, aber nicht um jeden Preis.
    In der Vergangenheit gab es dazu Anträge und Beschlüsse. Die sollte man
    wieder herausholen.

    Gunhild merkt an, dass es mal einen Ausschuss zu dem Semesterticket gab.

    Philipp sagt, dass es den Versuch von allen Hamburger ASten gab am Ball zu
    bleiben, dies aber auf der Strecke geblieben ist. Es gibt bereits Arbeit,
    an der angedockt werden kann.

    Artur merkt an, dass ein Begründungsteil erreicht wurde in dem Beitragsbescheid,
    der Infos zum Senatsprozent enthält.

    Mehrheitlich bei einer Gegenstimme und keiner Enthaltung angenommen.

    \subsection{Manifest zu G20}

    Artur stellt den Antrag vor.

    Yasemin fragt, warum gefordert wird die Tarife und Personalmittel im
    öffentlichen Dienst zu erhöhen, wo es durchaus Menschen gibt, die weit mehr
    als Mindestlohn verdienen.

    Tobias entgegnet, dass eine Erhöhung dort allen Beschäftigten zugute kommen
    und ein Verteilungskampf nicht gestartet werden sollte.

    22.11

    Franzi merkt an, dass Bezahlung im öffentlichen Dienst eher unter dem Einfluss
    von Aktiven steht als privatwirtschaftliche Verhältnisse. Es kann als Vorbild
    fungieren.

    Artur weist darauf hin, dass Personalmittel auch für mehr Einstellungen
    steht und dies sinnvoll ist, weil diese Gruppen in der Vergangenheit eine
    große Arbeitsverdichtung erfahren haben.

    Tobias fragt, warum Yasemin die geäußerte Position hat.

    Yasemin merkt an, dass der Satz unvermittelt wird.

    Tobias merkt an, dass die Stadt der größte Arbeitgeber in der Stadt ist.
    Die Bundesbank fordert mittlerweile bereits, dass Gewerkschaften echte
    Lohnforderungen stellen sollen. Ebenfalls weist er darauf hin, dass Deutschland
    ein Exportüberschuss von rund 200 Mrd. Euro hat und die in Lohnforderungen
    umgesetzt werden sollten.

    Franzi ...

    Timo hat nachgerechnet, dass für die Elbphilharmonie allen Kindern 2,5 Jahre
    Musikunterricht hätte gegeben werden können. Diese Menge Geld sollte daher
    jetzt nicht auch noch in den G20-Gipfel gesteckt werden.

    Franzi meint, dass Großteil der Bevölkerung Ahnung hat, wo Geld investiert
    werden könnte.

    Tobias stellt die Frage in den Raum, wie Staatswesen verfasst sein sollte.
    Prestige oder echte Arbeit an Verbesserung? Man müsste verstehen, dass
    der Topf Personalmittel einer Stadt wichtig ist, da wichtige Dinge Geld
    kosten und die Arbeitenden ihre Aufmerksamkeit auf sinnvolle Arbeit lenken
    können sollten und nicht auf ihren eigenen Haushalt.
    Tobias erwähnt ferner, dass es wert ist das Geld an anderer Stelle wegzunehmen.

    Franzi erwähnt, dass alle aufgerufen sind mit dem Manifest weitere Unterstützer
    zu suchen.

    Tobias fragt, wie man sich auf das Manifest beziehen kann.

    Franzi antwortet, dass das Manifest für Demomobilisierung verwendet werden soll.
    Pressekonferenz wird es dann im März geben.

    Mehrheitlich bei einer Gegenstimme und 3 Enthaltungen angenommen.

    22.29
    \subsection{VG-Wort}

    Geoffrey stellt den Antrag vor.

    Till wiederholt seinen Senf vom letzten Mal. Er hat noch immer Probleme
    mit dem Antrag. Er liest seine Alternativversion vor. Till schlägt vor
    für die nächste Sitzung einen Konsensvorschlag einzureichen.

    Karima findet Angriff auf Digitalisierung problematisch.

    Till findet, dass eine Forderung fehlt, dass in Bibliotheken genügend
    Materialien vorhanden sind.

    \subsection{Englischsprachige Wahlwerbung}

    Jim stellt seinen Antrag vor.

    Artur weist auf Mail vom "Präsidium" hin und erwähnt, dass es ein Problem
    ist, wenn Studierende hier nur zum Studieren herkommen. Außerdem sieht
    er ein Problem in der Übersetzungskapazität.

    Gunhild sieht ein Problem darin so alarmistisch englische Infos zu fordern.
    Sie meint, dass man hier Deutsch sprechen solle und auch noch genügend andere
    Dinge zu tun habe.

    Franzi sieht das Problem im Mangel an Finanzierung für Deutsch als Fremdsprache.

    Philipp sieht da kein Problem.

    Timo findet, dass es internationalen Studierenden nicht zugemutet werden kann,
    sich auf Deutsch zu informieren.

    Armin pflichtet Philipp bei. Geoffrey hätte kein Problem zumindest die
    allgemeinen Texte zu übersetzen.

    Till schauspielert und stellt einen Änderungsantrag vor.

    Es gibt eine Eskalation.

    Franzi hält ein Plädoyer für Deutsch als Fremdsprache und zeigt Probleme
    mit dem Studienkolleg und #uhhHilft auf. Die Studierenden müssten real in
    die Lage versetzt werden, die Sprache der Macht in Deutschland zu sprechen
    und sich an der politischen Diskussion zu beteiligen.
    In der Urnenwahlwoche sollte man auf dem Campus

    Melf stellt GO-Antrag zum Ende der Debatte. Alternative: Ende der Redeliste.

    8:10:2

    Die Redeliste wird abgearbeitet.

    Karima beschwert sich über die Metadebatte.

    Geoffrey stellt GO-Antrag die Sitzung nach dem TOP abzubrechen.

    M:0:1

    Laura beschwert sich über Redekultur. Sie empfiehlt Plakate aufzustellen,
    die informieren, dass z.B. Doktoranden und ERASMUS-Studierende wahlberechtigt
    sind. Außerdem sieht sie nicht den Konflikt der Diskussion.

    Mario sieht den erzieherischen Gedanken in der Bevormundung von Studierenden:
    Sie dürften nur wesentliche Infos in Englisch bekommen.

    Finanz-Timo findet man kann beides machen.

    Till sieht den Konflikt darin, ob man ein Problem einfach lösen kann oder
    ob das Problem nicht ausreichend reflektiert ist. Er findet alle Menschen
    in der Gesellschaft unmündig und will dies überwinden. Außerdem sieht er
    es als notwendig an, nicht mehr Arbeit dem Präsidium zuzuschanzen. Zweckdienliche
    Umsetzung.

    Sitzung um 23:42 abgebrochen.


    \subsection{Verschiedenes}
\end{document}
